\chapter{Évaluation Finale Globale}
\label{chap:evalfinale}

Cette évaluation couvre l'intégralité du programme : Gouvernance, Risque, Audit, Incident, Continuité et Crise.

% ------------------------------------------------------------------------------
% PARTIE A : QCM (Selection représentative)
% ------------------------------------------------------------------------------
\section{Partie A : QCM (Choisir la bonne réponse)}

\subsection*{Module 1 : Gouvernance et Risque}
\begin{enumerate}
    \item Le "Accountable" dans une matrice RACI est :
    \begin{itemize} \item[A.] Celui qui fait le travail. \item[B.] Celui qui rend des comptes. \item[C.] Celui qu'on informe. \end{itemize}
    \item La méthode d'analyse de risque française est :
    \begin{itemize} \item[A.] NIST. \item[B.] EBIOS RM. \item[C.] OCTAVE. \end{itemize}
    \item La criticité se calcule par :
    \begin{itemize} \item[A.] Probabilité + Gravité. \item[B.] Probabilité $\times$ Gravité. \item[C.] Coût $\times$ Temps. \end{itemize}
    \item Une vulnérabilité est :
    \begin{itemize} \item[A.] Une faiblesse. \item[B.] Un pirate. \item[C.] Un impact. \end{itemize}
    \item Souscrire une assurance cyber est une stratégie de :
    \begin{itemize} \item[A.] Réduction. \item[B.] Acceptation. \item[C.] Transfert. \end{itemize}
    \item Le RSSI reporte généralement à :
    \begin{itemize} \item[A.] La DSI. \item[B.] La Direction Générale. \item[C.] Le Service RH. \end{itemize}
    \item Le risque résiduel est :
    \begin{itemize} \item[A.] Le risque avant traitement. \item[B.] Le risque après traitement. \item[C.] Un risque imaginaire. \end{itemize}
    \item L'ISO 31000 concerne :
    \begin{itemize} \item[A.] La qualité. \item[B.] Le management du risque. \item[C.] L'audit. \end{itemize}
    \item L'appétence au risque est définie par :
    \begin{itemize} \item[A.] Le RSSI. \item[B.] La Direction (COMEX). \item[C.] L'auditeur. \end{itemize}
    \item Un actif essentiel est :
    \begin{itemize} \item[A.] Un serveur. \item[B.] Un processus métier critique. \item[C.] Un logiciel office. \end{itemize}
\end{enumerate}

\subsection*{Module 2 : Audit et Qualité}
\begin{enumerate}
    \setcounter{enumi}{10}
    \item L'indépendance de l'auditeur signifie :
    \begin{itemize} \item[A.] Travailler seul. \item[B.] Ne pas être juge et partie. \item[C.] Être externe. \end{itemize}
    \item La norme pour conduire un audit est :
    \begin{itemize} \item[A.] ISO 27001. \item[B.] ISO 19011. \item[C.] ISO 9001. \end{itemize}
    \item Une Non-Conformité Majeure implique :
    \begin{itemize} \item[A.] Un petit retard. \item[B.] Une défaillance systémique. \item[C.] Une suggestion. \end{itemize}
    \item La SOA (Statement of Applicability) est liée à :
    \begin{itemize} \item[A.] ISO 9001. \item[B.] ISO 27001. \item[C.] ISO 22301. \end{itemize}
    \item La correction sert à :
    \begin{itemize} \item[A.] Éliminer la cause. \item[B.] Éliminer le symptôme immédiat. \item[C.] Former le personnel. \end{itemize}
    \item L'audit de certification est réalisé par :
    \begin{itemize} \item[A.] L'entreprise elle-même. \item[B.] Un organisme certificateur. \item[C.] Un client. \end{itemize}
    \item COBIT est un référentiel pour :
    \begin{itemize} \item[A.] La gouvernance IT. \item[B.] La sécurité physique. \item[C.] La comptabilité. \end{itemize}
    \item La réunion d'ouverture sert à :
    \begin{itemize} \item[A.] Présenter les résultats. \item[B.] Confirmer le plan d'audit. \item[C.] Signer le certificat. \end{itemize}
    \item Une preuve d'audit peut être :
    \begin{itemize} \item[A.] Une rumeur. \item[B.] Un log ou un document. \item[C.] Une opinion. \end{itemize}
    \item Le cycle PDCA signifie :
    \begin{itemize} \item[A.] Plan, Do, Check, Act. \item[B.] Plan, Detect, Correct, Audit. \end{itemize}
\end{enumerate}

\subsection*{Module 3 : Incident et Continuité}
\begin{enumerate}
    \setcounter{enumi}{20}
    \item Un incident est :
    \begin{itemize} \item[A.] Un événement normal. \item[B.] Un événement compromettant la sécurité. \item[C.] Une crise médiatique. \end{itemize}
    \item L'ordre de volatilité recommande de collecter d'abord :
    \begin{itemize} \item[A.] Le disque dur. \item[B.] La RAM. \item[C.] Les sauvegardes. \end{itemize}
    \item Le RTO mesure :
    \begin{itemize} \item[A.] La perte de données. \item[B.] Le temps de reprise. \item[C.] Le coût. \end{itemize}
    \item Le RPO mesure :
    \begin{itemize} \item[A.] La perte de données acceptable. \item[B.] Le temps de panne. \item[C.] La gravité. \end{itemize}
    \item Le SIEM sert à :
    \begin{itemize} \item[A.] Surveiller et corréler les logs. \item[B.] Chiffrer les données. \item[C.] Gérer les projets. \end{itemize}
    \item Le confinement consiste à :
    \begin{itemize} \item[A.] Supprimer le virus. \item[B.] Limiter la propagation. \item[C.] Restaurer. \end{itemize}
    \item Un "Hot Site" est un site de secours :
    \begin{itemize} \item[A.] Vide. \item[B.] Allumé et synchronisé. \item[C.] Éteint. \end{itemize}
    \item Le BIA permet d'identifier :
    \begin{itemize} \item[A.] Les virus. \item[B.] Les fonctions vitales. \item[C.] Les failles. \end{itemize}
    \item ISO 27035 traite de :
    \begin{itemize} \item[A.] La qualité. \item[B.] Les incidents de sécurité. \item[C.] La cryptographie. \end{itemize}
    \item Un "Faux Positif" est :
    \begin{itemize} \item[A.] Une alerte injustifiée. \item[B.] Un virus manqué. \item[C.] Une erreur humaine. \end{itemize}
\end{enumerate}

\subsection*{Module 4 : Crise et Retex}
\begin{enumerate}
    \setcounter{enumi}{30}
    \item Une crise se caractérise par :
    \begin{itemize} \item[A.] Une panne. \item[B.] Une menace pour la survie. \item[C.] Un bug. \end{itemize}
    \item Le rôle du Scribe est de :
    \begin{itemize} \item[A.] Parler à la presse. \item[B.] Noter la chronologie. \item[C.] Restaurer le SI. \end{itemize}
    \item La règle des 3T inclut :
    \begin{itemize} \item[A.] Technique, Test, Temps. \item[B.] Ténor, Transparence, Temporalité. \item[C.] Travail, Tâche, Terme. \end{itemize}
    \item Le RETEX à froid se fait :
    \begin{itemize} \item[A.] Immédiatement. \item[B.] Quelques semaines après. \item[C.] Jamais. \end{itemize}
    \item La méthode des 5 Pourquois sert à trouver :
    \begin{itemize} \item[A.] Le coupable. \item[B.] La cause racine. \item[C.] Le coût. \end{itemize}
    \item Le Porte-parole doit être :
    \begin{itemize} \item[A.] Unique. \item[B.] Multiple. \item[C.] Technique. \end{itemize}
    \item Un "Holding Statement" est :
    \begin{itemize} \item[A.] Un rapport final. \item[B.] Une déclaration initiale. \item[C.] Un plan technique. \end{itemize}
    \item La cellule de crise est composée de :
    \begin{itemize} \item[A.] Techniciens uniquement. \item[B.] Décideurs et Communicants. \item[C.] Clients. \end{itemize}
    \item En communication de crise, il faut :
    \begin{itemize} \item[A.] Mentir pour protéger. \item[B.] Être transparent. \item[C.] Ne rien dire. \end{itemize}
    \item Le PGC est :
    \begin{itemize} \item[A.] Plan de Gestion de Crise. \item[B.] Plan de Gestion Commerciale. \end{itemize}
\end{enumerate}

% ------------------------------------------------------------------------------
% PARTIE B : VRAI / FAUX (Selection représentative)
% ------------------------------------------------------------------------------
\section{Partie B : Vrai ou Faux}

\textbf{Module 1 :}
\begin{enumerate}
    \item Le DSI est le seul responsable de la sécurité (Accountable). \textbf{(Faux - C'est le RSSI/PDG)}
    \item Le risque zéro n'existe pas. \textbf{(Vrai)}
    \item EBIOS RM est une méthode américaine. \textbf{(Faux - Française ANSSI)}
    \item Le transfert de risque consiste à accepter la perte. \textbf{(Faux - C'est l'assurance ou externalisation)}
    \item Un actif peut être intangible (réputation). \textbf{(Vrai)}
\end{enumerate}

\textbf{Module 2 :}
\begin{enumerate}
    \setcounter{enumi}{5}
    \item On peut auditer son propre travail. \textbf{(Faux - Interdit)}
    \item ISO 19011 donne les exigences du SMSI. \textbf{(Faux - C'est ISO 27001)}
    \item Une observation est une non-conformité majeure. \textbf{(Faux - C'est une piste d'amélioration)}
    \item L'audit interne prépare souvent la certification. \textbf{(Vrai)}
    \item Le rapport d'audit ne contient que des opinions. \textbf{(Faux - Que des faits/preuves)}
\end{enumerate}

\textbf{Module 3 :}
\begin{enumerate}
    \setcounter{enumi}{10}
    \item Il faut éteindre immédiatement un serveur piraté. \textbf{(Faux - Risque de perte de preuves RAM)}
    \item Le RTO est toujours plus long que le RPO. \textbf{(Vrai généralement, sauf pertes totales)}
    \item Le PCA est un plan purement informatique. \textbf{(Faux - C'est un plan Métier)}
    \item Un site froid est opérationnel immédiatement. \textbf{(Faux - C'est le site chaud)}
    \item Le Hash garantit l'intégrité d'un fichier. \textbf{(Vrai)}
\end{enumerate}

\textbf{Module 4 :}
\begin{enumerate}
    \setcounter{enumi}{15}
    \item Le temps médiatique est plus lent que le temps technique. \textbf{(Faux - Plus rapide)}
    \item Le RETEX sert à chercher les coupables. \textbf{(Faux - A améliorer le système)}
    \item On peut avoir plusieurs porte-paroles en même temps. \textbf{(Faux - Une seule voix)}
    \item La transparence absolue est obligatoire, même sur les négociations de rançon. \textbf{(Faux - Secrét sur les rançons)}
    \item L'exercice Tabletop est une simulation physique. \textbf{(Faux - C'est un exercice sur table/débat)}
\end{enumerate}

% ------------------------------------------------------------------------------
% PARTIE C : Questions Ouvertes (Selection représentative)
% ------------------------------------------------------------------------------
\section{Partie C : Questions Ouvertes}

\begin{enumerate}
    \item \textbf{(Risque)} Expliquez la différence entre "Vulnérabilité" et "Menace" avec un exemple.
    \item \textbf{(Gouvernance)} Pourquoi le RSSI ne doit-il pas être sous la responsabilité hiérarchique directe de la DSI ?
    \item \textbf{(Audit)} Décrivez les 3 sources de preuves possibles lors d'un audit.
    \item \textbf{(Audit)} Quelle est la différence entre une "Correction" et une "Action Corrective" ?
    \item \textbf{(Incident)} Expliquez l'ordre de volatilité des données. Pourquoi est-il important ?
    \item \textbf{(Continuité)} Définissez le RTO et le RPO. Quel lien avec le coût de la solution ?
    \item \textbf{(Crise)} Quels sont les 3 rôles principaux au sein d'une cellule de crise ?
    \item \textbf{(Communication)} Expliquez la règle des "3 T" en communication de crise.
    \item \textbf{(Retex)} Quelle est l'utilité de la méthode des "5 Pourquois" ?
    \item \textbf{(Synthèse)} Comment le cycle PDCA est-il illustré tout au long de ce cours (du risque au RETEX) ?
\end{enumerate}